\section{Exact verification of the algebraic certificate}
\label{sec:certificate-proof}

We return to the free algebraic data of Section~\ref{sec:certificate}
to prove Proposition~\ref{cert:universal}. All tensors are again over
$\R^5$. The scalars $w,z$ and vector $\ell$ are free, and $X$ is an
arbitrary symmetric trace-free three-tensor.
Table~I specifies the polynomials $f,U,K,M,P$, and Table~II specifies
nonnegative comparison forms. For the interior inequality, we compare
$\mathcal I$ with a positive diagonal quadratic form in $(w,z,\ell,X)$
plus projected tensor squares, and bound the remaining error term by term.
For the boundary inequality, the analogous comparison uses scalar squares
in the four tangential eigenvalues and a uniform bound on eleven remainder
coefficients. We give the contraction rules and rational bounds in full.

\subsection{Tensor monomials and the fixed polynomials}\label{cert:decode}

Here and in Appendix~\ref{app:tables}, a rank-$r$ tensor monomial is encoded by
a word of decimal digits, partitioned into triples.  Its first triple $hbj$
specifies
\[
 g^h(\tr A^3)^b w^{\mathbf1_{\{j=1\}}}
                         z^{\mathbf1_{\{j=2\}}}.
\]
Each later triple $uvk$ specifies an edge carrying the matrix $A^k$, with
$A^0=\delta$.  An endpoint labeled $i\in\{0,\ldots,r-1\}$ is the $i$th free
slot, starting with slot zero.  An endpoint labeled $5$ carries one copy of
$p$, and one labeled $6$ carries one copy of $\ell$.  Three incidences of label
$7$ are contracted against the three slots of one $X$; three incidences of
label $8$, if present, use a second copy of $X$.  Each occurrence of a vector
endpoint is a separate vector factor.  Matrix entries, vector entries, and
tensor entries are multiplied, and all nonfree slots are summed.  Edges are
unoriented because $A$ is symmetric.  Permuting the three incidences at an
$X$ does not change its value because $X$ is symmetric.  There are no
implicit multiplicities or factorials in this convention.

Write $\mathcal M_r(s)$ for the tensor encoded by the word $s$.  For example,
\[
 \mathcal M_0(\texttt{200550})=g^2|p|^2,\qquad
 \mathcal M_4(\texttt{000500520131})_{lijk}=p_lp_jA_{ik},\qquad
 \mathcal M_1(\texttt{000771700})_l=q_l.
\]
The digit labels encode slots and must be distinguished from the coordinate
indices that are summed from $1$ to $5$.

\subsubsection{Table I}\label{cert:table-I}

For $F\in\{f,U,K,M,P\}$, let $r_F$ be respectively $0,2,4,1,1$, with slot
orders respectively empty, $(j,k)$, $(l,i,j,k)$, $(l)$, $(l)$.  If a row of
Table~I consists of $F,s,c$, its contribution is $c\mathcal M_{r_F}(s)/Q$;
thus each polynomial is given by
\begin{equation}\label{cert:fixed-polynomials}
 F=\frac1Q\sum_{(F,s,c)\ \text{in Table I}}c\,\mathcal M_{r_F}(s).
\end{equation}
Table~I contains $34,57,34,74,32$ rows for $f,U,K,M,P$, respectively.
Every free slot occurs exactly once in every row, and these rows have no
$X,\ell,w,z$ factors.  Every $P$ row has a positive $g$ exponent.  Because
an edge has two endpoints and $U,K$ have even rank, their numbers of $p$
endpoints are even.  This proves the parity assertions of
Proposition~\ref{cert:universal} directly from the definitions.  The same
definitions show that all five polynomials are equivariant under orthogonal
changes of coordinates, with $f$ invariant.

\subsection{The nonnegative comparison polynomials}\label{cert:gram}

The positive tensor squares below have the form of Gram-matrix
certificates for polynomial positivity; compare
\cite[Section~3.2]{Parrilo2003}. Here all rational polynomials,
projection identities, and remainder bounds are explicit, and we verify
the certificate directly.


A header $\texttt{@ }r\ \mathrm{kind}\ \mathrm{mode}$ in Table~II begins a
new block; consecutive equal headers still designate different blocks.
For a block of nonnegative rank, write its monomial rows as $s_1,\ldots,s_N$
and its rectangular integer coefficient matrix as $(c_{tc})$.  Define
\[
 J_c=\frac1Q\sum_{t=1}^N c_{tc}\mathcal M_r(s_t).
\]
Its contribution to $\Gamma$ is
\begin{equation}\label{cert:gram-block}
 \rho\sum_c\langle J_c,\Pi J_c\rangle,\qquad
 \rho=\begin{cases}1,&\mathrm{mode}=\texttt{I},\\
                    \chi,&\mathrm{mode}=\texttt{B}.
       \end{cases}
\end{equation}
Here $\Pi$ is specified below by $r$ and $\mathrm{kind}$.  There are nineteen
such blocks with a total of $529$ monomial rows.  Every row has degree one in
$(w,z,\ell,X)$, assigning degree one to each component of these variables and
degree zero to $A,p,g$.  Hence $\Gamma$ has degree two in these variables.

The two remaining blocks have rank $-1$ and fourteen rows each.  In these
blocks the four digits of a row word are an exponent
$\beta=(\beta_1,\ldots,\beta_4)$, encoding $x^\beta$ for $x\in\R^4$.
Both coefficient matrices have five columns.  Set
\[
 J_c(x)=\frac1Q\sum_t c_{tc}x^{\beta_t},\qquad
 \operatorname{Av}F(x)=\frac1{24}\sum_{\pi\in S_4}
                  F(x_{\pi(1)},\ldots,x_{\pi(4)}).
\]
Their sum defines the scalar polynomial
\begin{equation}\label{cert:boundary-gram}
 \Lambda(x)=\sum_{\text{rank }-1\text{ blocks}}
            \operatorname{Av}\left(\rho(x)\sum_c J_c(x)^2\right),
 \qquad
 \rho(x)=\begin{cases}1,&\mathrm{mode}=\texttt{I},\\
             x_1+x_2+x_3+x_4,&\mathrm{mode}=\texttt{H}.
             \end{cases}
\end{equation}
Consequently $\Lambda\ge0$ when $\sum x_i\ge0$.

\subsubsection{Scaled projection operators}\label{cert:projections}

For an arbitrary rank-$r$ tensor $F$, let $F^\pi$ denote slot permutation
by $\pi$.  We use the averaged symmetrization
$\Sym F=(r!)^{-1}\sum_{\pi\in S_r}F^\pi$.
The operator $\operatorname{Tr}_{ab}$ contracts slots $a,b$ with the
Euclidean metric, leaving the remaining slots in their original order.
The label $\texttt{none}$ means $\Pi=\Id$.  The label
$\texttt{anti}$ means
\[
 \Pi F=\sum_{\pi\in S_r}\operatorname{sgn}(\pi)F^\pi.
\]
For $\texttt{sym}$, put $F'=\sum_{\pi\in S_r}F^\pi$ and
$v=\operatorname{Tr}_{12}F'$.  The operators used in Table~II are
\begin{align}
 (\Pi F)_{ij}&=5F'_{ij}-\delta_{ij}v &&(r=2),\label{cert:sym2}\\
 (\Pi F)_{ijk}&=7F'_{ijk}
       -(\delta_{ij}v_k+\delta_{ik}v_j+\delta_{jk}v_i) &&(r=3).
       \label{cert:sym3}
\end{align}
For $r=3$ and $\texttt{hook}$, set
\begin{align}
 F'&=3(F+F^{(12)})-\sum_{\pi\in S_3}F^\pi,
       &v&=\operatorname{Tr}_{12}F',\notag\\
 (\Pi F)_{ijk}&=8F'_{ijk}
       -(2\delta_{ij}v_k-\delta_{ik}v_j-\delta_{jk}v_i).
       \label{cert:hook}
\end{align}

\begin{lemma}\label{cert:positive-projections}
All these operators are self-adjoint and positive semidefinite on their
entire tensor spaces.  In particular, $\Gamma\ge0$ whenever $\chi\ge0$.
\end{lemma}
\begin{proof}
Each slot permutation is orthogonal, with adjoint its inverse.  Averaging
over all permutations, with or without their signs, gives the orthogonal
projection onto symmetric or alternating tensors.  Thus $\texttt{anti}$ is
$r!$ times an orthogonal projection.  For a symmetric matrix $S$,
$S-\delta\tr S/5$ is its orthogonal trace-free projection.  Since $F'=2\Sym F$
in \eqref{cert:sym2}, that operator is ten times the symmetric trace-free
projection.

For rank three define
\[
 (\mathcal Jv)_{ijk}=\delta_{ij}v_k+\delta_{ik}v_j+\delta_{jk}v_i.
\]
Then $\operatorname{Tr}_{12}\mathcal Jv=7v$, and for symmetric $S$,
$\langle S,\mathcal Jv\rangle=3\langle\operatorname{Tr}_{12}S,v\rangle$.
It follows that $S-\mathcal J(\operatorname{Tr}_{12}S)/7$ is the orthogonal
projection onto symmetric trace-free tensors.  Since now $F'=6\Sym F$,
\eqref{cert:sym3} is $42$ times that projection.

For completeness, the trace correction in \eqref{cert:hook} can be checked
without any representation-theoretic assumptions.  Let
\[
 S_{12}=\frac{\Id+(12)}2,\qquad
 S=\frac16\sum_{\pi\in S_3}\pi,\qquad P_0=S_{12}-S.
\]
The identities $S_{12}S=SS_{12}=S$ show that $P_0$ is an orthogonal
projection.  Its image $\mathcal W$ consists of first-pair symmetric tensors
orthogonal to fully symmetric tensors.  For $H\in\mathcal W$,
\[
 H_{ijk}=H_{jik},\qquad H_{ijk}+H_{ikj}+H_{jki}=0.
\]
If $v=\operatorname{Tr}_{12}H$, contraction of the second identity gives
$\operatorname{Tr}_{13}H=\operatorname{Tr}_{23}H=-v/2$.
Define
\[
 (\mathcal Bv)_{ijk}=2\delta_{ij}v_k-\delta_{ik}v_j-\delta_{jk}v_i.
\]
Direct substitution gives $\mathcal Bv\in\mathcal W$,
$\operatorname{Tr}_{12}\mathcal Bv=8v$, and
\[
 \langle H,\mathcal Bz\rangle
 =2\langle v,z\rangle
  -\langle\operatorname{Tr}_{13}H,z\rangle
  -\langle\operatorname{Tr}_{23}H,z\rangle
 =3\langle v,z\rangle.
\]
Therefore $Q_0H=H-\mathcal B(\operatorname{Tr}_{12}H)/8$ has all traces
zero, and $H-Q_0H$ is orthogonal to every trace-free tensor in $\mathcal W$.
Thus $Q_0$ is the orthogonal projection in $\mathcal W$ onto its trace-free
part.  Since $F'=6P_0F$, the hook operator is $48Q_0P_0$; $Q_0P_0$ is the
orthogonal projection in the full tensor space onto that subspace.  In
particular
\[
 \langle F,\Pi F\rangle=48|Q_0P_0F|^2\ge0.
\]
Finally the weights in \eqref{cert:gram-block} are nonnegative under the stated
constraint.
\end{proof}

\subsection{Exact contraction and differentiation rules}\label{cert:collection}

We give explicit collection rules both to specify the finite calculation and
to justify that it is an identity on all the constrained tensors.  A graph
key records the rank, powers of $g,w,z$, a multiset of trace factors
$\tr A^k$, and a multiset of edges with their powers.  For purposes of
ordering keys, replace endpoint labels $5,6,7,8$ by $-1,-2,-3,-4$.
Sort the endpoints of each edge, sort the edge list and trace list, and take
the lexicographically larger edge list obtained by retaining or exchanging
$-3,-4$.  A lone cubic tensor is thereby labeled $-3$.  These choices are
only ordering conventions: exchange of two identical copies of $X$ preserves
the value of every graph.

Tensor multiplication takes the disjoint union of graphs, shifts the free
slot labels of the second factor, and distinguishes any two cubic tensor
copies.  A contraction of two free slots joins their incident edges.  If
these are different edges, their powers add; if they are the same edge, it
becomes a trace $\tr A^k$.  The removed free slots are deleted and the others
renumbered in order.  The identities underlying these operations are precisely
\[
 \sum_i(A^h)_{ai}(A^k)_{ib}=(A^{h+k})_{ab},\qquad
 \sum_i(A^k)_{ii}=\tr A^k.
\]
Reduce the trace factors of powers $0,1,2$ to $5,0,1$, respectively.  An
edge of power zero connecting two slots of the same $X$ makes the monomial
zero.  Thus every reduction is an equality on the allowed data.  Other
identities need not be imposed: retaining two different keys for functions
that happen to agree cannot invalidate either the expansion or the later
triangle inequality.

The derivative of a path and of a trace is
\begin{align}
 D_l(A^k)&=\sum_{h=0}^{k-1}A^hX_lA^{k-1-h}-kq_lA^k,
       & (X_l)_{ij}&=X_{lij},\label{cert:path-derivative}\\
 D_l\tr A^k&=k\,\tr(A^{k-1}X_l)-kq_l\tr A^k.
       \label{cert:trace-derivative}
\end{align}
For $k=0$ the sums and derivatives are zero.  These formulas follow by
differentiating each matrix factor using \eqref{cert:D}; cyclic invariance
of trace makes all $k$ insertion positions on a trace equal.  In graph form,
inserting $X_l$ on an edge of power $k$ splits it into two paths of lengths
$h,k-1-h$ ending at a new $X$, whose last slot is the derivative index $l$.
The negative terms from all $A$ factors combine into minus the total
$A$ degree times $q_l$ times the original graph.  Each $p$ endpoint is in
turn replaced by $wA_{l\cdot}-gz\delta_{l\cdot}$, and a factor $g^h$ gives
$hzg^{h-1}p_l$.  If an edge has two $p$ endpoints, these are two product-rule
terms, including their multiplicity when they coincide.

These operations are also sound if traces were simplified first, since
\[
 D_l\tr A=0,\qquad
 D_l\tr A^2=2X_{lij}A_{ij}-2q_l|A|^2=0.
\]
Only the original Table~I monomials are differentiated, exactly once.  No
already differentiated graph, or graph involving $X,w,z$, is differentiated
again.  Thus no unspecified higher derivative enters the calculation.

\begin{lemma}[Closure of the graph calculation]\label{cert:closure}
The preceding rules evaluate all terms of $\mathcal I$ and $\Gamma$ by a
finite sum of scalar graphs.  In every intermediate graph a free slot has
exactly one incidence, each cubic tensor has exactly three incidences, and
there are at most two cubic tensors.  Every resulting scalar graph has
degree two in $(w,z,\ell,X)$.
\end{lemma}
\begin{proof}
The incidence assertions hold for each table row by its decoding.  A
disjoint union preserves them.  A permutation only renames free slots.
Contracting two free slots removes their unique incidences and joins the
remaining ends, or closes a trace; it does not change the number of slots of
any cubic tensor.  The derivative insertion described above replaces one
edge by two incidences of a new $X$ and supplies its third incidence as a
new free slot.  Trace differentiation does the same with a self-edge.
The $q_l$ term supplies two incidences on its self-edge and the third at $l$.
Derivatives of $p$ and $g$ supply a free derivative slot without introducing
a cubic tensor.  These statements prove the incidence assertions by
induction over the operations.

Table~I has degree zero and its derivative has degree one in the indicated
variables.  The formula for $\mathcal I$ multiplies such derivatives by one
further factor, or a degree-zero coefficient by two factors.  Each Table~II
row has degree one; projection, contraction, and multiplication by $\chi$
preserve that degree.  Squaring produces degree two.  In particular no
operation used here needs three copies of $X$.  Termination is immediate
from the finite tables, finite permutation sums, and the finite product
rule; every contraction reduces the number of free slots by two.
\end{proof}

The linear contribution of each Table~I row is obtained by substitution into
\eqref{cert:interior}, forming $C=K-K^{(12)}$ for a $K$ row.  One check on
the slot order and coefficients is the single row $U=A$ with coefficient
one.  Its interior contribution is
\[
 5z^2+|X|^2+(2a-2)|q|^2,
\]
and its boundary contribution is $2\lambda-\tr A^3$.
To expand a Gram block, sum over unordered pairs of rows $s,t$: the
coefficient is $\sum_c c_{sc}c_{tc}/Q^2$, doubled when $s\ne t$, and the
graph polynomial is the contraction of row $t$ with $\Pi$ applied to row
$s$, multiplied by $\rho$.  This uses self-adjointness of $\Pi$ and accounts
for every term of \eqref{cert:gram-block} exactly once.

\subsection{Bounds for scalar graphs}\label{cert:bounds}

Put $\theta=4/5$.  We first establish the tensor estimates that will bound
each scalar graph, retaining the dependence between $p$ and $g$.

\begin{lemma}\label{cert:tensor-bounds}
For the constrained matrix $A$ and symmetric trace-free tensor $X$,
\begin{align}
 \|A\|_{\mathrm{op}}^2&\le\theta,&
 |A^k|^2&\le\theta^{k-1}\quad(k\ge1),\label{cert:matrix-bounds}\\
 |\tr A^k|^2&\le\theta^{k-2}\quad(k\ge2),&
 |\tr A^3|^2&\le\frac9{20},\label{cert:trace-bounds}\\
 |X:A|^2&\le\frac57|X|^2,&
 (X:A)_i&=X_{ijk}A_{jk}.\label{cert:five-sevenths}
\end{align}
If $h,b>0$ and $g^2+|p|^2\le1$, then
\begin{equation}\label{cert:mixed-bound}
 |g|^{2h}|p|^{2b}
 \le \left(\frac{h}{h+b}\right)^h
       \left(\frac{b}{h+b}\right)^b.
\end{equation}
\end{lemma}
\begin{proof}
Let the eigenvalues of $A$ be $\alpha_1,\ldots,\alpha_5$.  For each $i$,
\[
 \alpha_i^2=\left(\sum_{j\ne i}\alpha_j\right)^2
       \le4\sum_{j\ne i}\alpha_j^2=4(1-\alpha_i^2).
\]
This proves the operator bound, and
$\sum_i\alpha_i^{2k}\le\theta^{k-1}\sum_i\alpha_i^2$ proves the
Frobenius bound.  Moreover
$|\sum_i\alpha_i^k|\le\sum_i|\alpha_i|^k
\le\theta^{(k-2)/2}\sum_i\alpha_i^2$ for $k\ge2$.

For the sharper cubic bound maximize and minimize $\sum_i\alpha_i^3$ on
the compact set $\sum_i\alpha_i=0$, $\sum_i\alpha_i^2=1$.  The gradients
of the two constraints are independent there.  Lagrange multipliers give
$3\alpha_i^2=\mu+2\nu\alpha_i$, so an extremum has two distinct
eigenvalues (a single value would violate the constraints).  If their
multiplicities are $j,5-j$, the values are, up to simultaneous sign,
\[
 \sqrt{\frac{5-j}{5j}},\qquad
 -\sqrt{\frac{j}{5(5-j)}}.
\]
Their cubic trace has absolute value
$|5-2j|/\sqrt{5j(5-j)}$.  Its square is $9/20$ for $j=1,4$ and $1/30$
for $j=2,3$, proving \eqref{cert:trace-bounds}.

To prove \eqref{cert:five-sevenths}, fix a vector $\zeta$ and symmetrize
$\zeta\otimes A$:
\[
 B_{ijk}=\frac{\zeta_iA_{jk}+\zeta_jA_{ik}+\zeta_kA_{ij}}3,
 \qquad v=\operatorname{Tr}_{12}B=\frac23 A\zeta.
\]
Expanding the nine terms in $|B|^2$ gives
$|B|^2=|\zeta|^2/3+2|A\zeta|^2/3$.
The symmetric trace-free projection is $F=B-\mathcal Jv/7$, in the notation
of Lemma~\ref{cert:positive-projections}.  Since
$|\mathcal Jv|^2=21|v|^2$ and
$\langle B,\mathcal Jv\rangle=3|v|^2$,
\[
 |F|^2=|B|^2-\frac37|v|^2
       =\frac13|\zeta|^2+\frac{10}{21}|A\zeta|^2
       \le\frac57|\zeta|^2.
\]
Symmetry and trace-freeness of $X$ give
$\langle X:A,\zeta\rangle=\langle X,F\rangle$; Cauchy--Schwarz and
duality prove the claim.

Finally put $s=g^2,t=|p|^2$.  For $h,b>0$, the maximum of $s^ht^b$ on
$s,t\ge0$, $s+t\le1$ lies on $s+t=1$.  Differentiating
$h\log s+b\log(1-s)$ gives the unique interior maximum
$s=h/(h+b)$, proving \eqref{cert:mixed-bound}.  If either exponent is zero,
the corresponding bound is simply $1$.
\end{proof}

For a nonzero scalar graph key $t$ of degree two, let
\[
 y=(|w|,|z|,|\ell|,|X|)^T
\]
and let $i,j\in\{1,2,3,4\}$ be its two degree-one types, counted with
multiplicity and ordered so $i\le j$.  Let $h$ be its $g$ exponent and $b$
its number of $p$ endpoints.  Define $L_t^2$ as a product of the following
rational factors:
\begin{enumerate}
\item\label{cert:graph-rule-trace}
      for each trace $\tr A^k$, use $\theta^{k-2}$, replacing this by
      $9/20$ when $k=3$;
\item\label{cert:graph-rule-edge}
      for an edge of power $k$ between distinct vertices, use $\theta^k$;
\item\label{cert:graph-rule-self}
      for an edge between two slots of the same $X$, use $5/7$ when $k=1$
      and $\theta^{k-1}$ when $k\ge2$;
\item\label{cert:graph-rule-mixed}
      if $h,b>0$, include the right side of \eqref{cert:mixed-bound};
      otherwise include $1$.
\end{enumerate}
Here separate occurrences of a vector endpoint count as distinct vertices;
thus a $p$--$p$ edge is governed by Item~\ref{cert:graph-rule-edge}, not
Item~\ref{cert:graph-rule-self}.  Power-zero self-edges
at $X$ have already vanished in the collection rules.

\begin{lemma}[Complete graph bound]\label{cert:graph-bound}
Every scalar graph retained by Lemma~\ref{cert:closure} satisfies
\begin{equation}\label{cert:monomial-bound}
 |\mathcal M(t)|\le L_t y_i y_j
\end{equation}
on $g^2+|p|^2\le1$.
\end{lemma}
\begin{proof}
First separate all scalar trace factors and use Lemma~\ref{cert:tensor-bounds}.
A cubic tensor has at most one self-edge, since such an edge consumes two
of its three incidences.  An edge of power $k\ge2$ reduces that tensor to
the vector $v_i=X_{ijk}(A^k)_{jk}$, for which
$|v|\le|X|\,|A^k|\le|X|\theta^{(k-1)/2}$.  For $k=1$ use
\eqref{cert:five-sevenths}.  These reductions leave at most two cubic
vertices and no self-edge on either.

We list all possible connected components of the remaining graph.
With zero cubic vertices, a component is one edge joining two vectors;
its contraction is bounded by the operator norm of its matrix times the
two vector norms.  With one cubic vertex, all its three slots end at vector
leaves.  Move each edge matrix onto its vector, and use tensor
Cauchy--Schwarz, obtaining the norm of $X$ times the three vector norms
and the product of the three operator norms.

If a component contains two cubic vertices, they are joined by $r=1,2,$
or $3$ edges; every other slot ends at a vector leaf.  Contract these leaves
first, obtaining rank-$r$ tensors $Y,Z$ with the same preceding norm
bounds.  Up to a permutation of their slots, the remaining contraction is
\[
 \left\langle Y,
      (A^{k_1}\otimes\cdots\otimes A^{k_r})Z\right\rangle.
\]
Its absolute value is at most
$|Y||Z|\prod_{s=1}^r\|A^{k_s}\|_{\mathrm{op}}$.  This includes all three
parallel edges when $r=3$, without a dimension factor.  If the two cubic
vertices have no edge between them, they lie in two components of the
one-vertex type already treated.  These cases are exhaustive: all other
vertices have valence one, so cannot occur internally on a path, and there
are no further cubic vertices.  The reduced vectors from self-edges are
ordinary vector leaves in this classification.

Multiply the estimates over components.  Each original $\ell$ or $X$
norm appears once per occurrence, all ordinary edge matrices contribute
$\theta^{k/2}$, and the $p$ endpoints contribute $|p|^b$.  The scalar
factors $g^h$ and $|p|^b$ are bounded jointly by
\eqref{cert:mixed-bound}.  Finally the total degree is two, so all remaining
$w,z,\ell,X$ factors give exactly $y_i y_j$.
\end{proof}

We can now bound every term left after subtracting the nonnegative
comparison forms. The graph estimates reduce this remainder to a
quadratic form in four magnitudes. The next calculation shows that
its error is smaller than the reserved positive diagonal terms.

\subsection{The interior remainder and its coercive margin}\label{cert:interior-remainder}

Set
\begin{equation}\label{cert:remainders}
 d=\left(\frac2{25},\frac2{25},\frac14,1\right),\qquad
 E=\Gamma+\sum_{i=1}^4d_i y_i^2-\mathcal I.
\end{equation}
The squared norms in Equation~\eqref{cert:remainders} are polynomial
contractions.  Expand
and collect $E=\sum_t e_t\mathcal M(t)$ by Section~\ref{cert:collection},
discarding zero coefficients only after their signed contributions have
been summed.  All coefficients are rational.  The denominator $Q^2$
suffices: Gram terms have that denominator, and the denominators of
$a,d_i$, and all Table~I coefficients divide it after multiplication by
$Q^2$.

For each retained key let $n_t$ be the least nonnegative integer satisfying
\begin{equation}\label{cert:root-rounding}
 n_t^2\ge10^{10}L_t^2,
 \qquad \widehat L_t=\frac{n_t}{10^5}.
\end{equation}
There is no floating-point square root in this prescription.  For a rational
$s=10^{10}L_t^2$, compute $n=\lfloor\sqrt{\lfloor s\rfloor}\rfloor$ by
integer arithmetic and increase it by one exactly when $n^2<s$.
Form a symmetric matrix $D_0$ by adding
\begin{equation}\label{cert:D0}
 \frac{|e_t|\widehat L_t}{1+\mathbf1_{\{i\ne j\}}}
\end{equation}
to its $ij$ entry and also its $ji$ entry if $i\ne j$, starting with zero.
The factor $1/2$ in an off-diagonal entry is required because the associated
term in $y^TD_0y$ is $2(D_0)_{ij}y_iy_j$.  Lemma~\ref{cert:graph-bound}
and the triangle inequality therefore give
\begin{equation}\label{cert:error-bound}
 |E|\le y^TD_0y.
\end{equation}

For clarity, the exact finite calculation needed here is the following.
Decode all rows of Tables~I and II; expand the nineteen Gram blocks by
unordered row pairs and the finite projections
\eqref{cert:sym2}--\eqref{cert:hook}; subtract each Table~I contribution
to \eqref{cert:interior}; and add the four contractions with coefficients
$d_i$.  Use the contraction, derivative, and key ordering rules above after
each operation.  Combine all signed coefficients at equal keys.  For each
nonzero key compute its rational $L_t^2$, perform
\eqref{cert:root-rounding}, and add \eqref{cert:D0}.
This gives all ten upper-triangular entries as follows:
\begin{equation}\label{cert:exact-interior-sums}
\begin{array}{c|r|r}
 (i,j)&\text{nonzero keys}&2\cdot10^5Q^2(D_0)_{ij}\\\hline
 (1,1)&89&321268243322644\\
 (1,2)&53&96616799870252\\
 (1,3)&123&294707669526536\\
 (1,4)&170&385792585960684\\
 (2,2)&64&269583307398596\\
 (2,3)&85&115717471031956\\
 (2,4)&106&125155603107368\\
 (3,3)&238&605509501361296\\
 (3,4)&414&618424270356936\\
 (4,4)&372&1004504025266456
\end{array}
\end{equation}
In particular there are $1714$ nonzero keys.  These finite sums use only
integer and rational operations and comparisons; the complete literal input
is in Appendix~\ref{app:tables}.  The accompanying program
\texttt{verification/checks/reproduce.py} implements the specified calculation and records
the individual coefficients as well as these sums.

Entrywise, \eqref{cert:exact-interior-sums} implies $D_0\le\overline D$, where
\begin{equation}\label{cert:upper-matrix}
 \overline D=\frac1{10^5}
 \begin{pmatrix}
 2571&773&2358&3087\\
 773&2157&926&1002\\
 2358&926&4845&4948\\
 3087&1002&4948&8037
 \end{pmatrix}.
\end{equation}
For example, this entrywise check consists of comparing each last-column
integer in \eqref{cert:exact-interior-sums} with the corresponding integer
in \eqref{cert:upper-matrix} multiplied by $2Q^2$.
Let $\nu=(2484,1440,996,304)^T$.  The exact diagonal slacks are
\begin{equation}\label{cert:slacks}
 \left(d_i-\frac{(\overline D\nu)_i}{\nu_i}\right)_{i=1}^4
 =\left(\frac{673}{18400},\frac{438907}{12000000},
         \frac{2844869}{24900000},\frac{3479389}{7600000}\right).
\end{equation}
For $y_i\ge0$, weighted Young inequalities give
\[
 2y_iy_j\le\frac{\nu_j}{\nu_i}y_i^2+
                  \frac{\nu_i}{\nu_j}y_j^2,
 \qquad
 y^TD_0y\le y^T\overline D y
       \le\sum_i\frac{(\overline D\nu)_i}{\nu_i}y_i^2.
\]
Use $\Gamma\ge0$ in \eqref{cert:remainders} and then
\eqref{cert:error-bound}.  It follows that $\mathcal I$ is at least the
sum of $y_i^2$ times the four positive slacks in \eqref{cert:slacks}.
In particular the coefficients of $z^2$ and $|\ell|^2$ are at least
$\epsilon_*$, proving \eqref{cert:interior}.  To obtain the stated positivity
of $f$, keep $A,p,g$ fixed, set $X=w=z=0$, and choose $|\ell|=1$.
Then $q=0$, every $D$ term is zero, and $\mathcal I=f$.

\subsection{Boundary reduction and the eleven coefficients}\label{cert:boundary-remainder}

It remains to verify \eqref{cert:boundary}.  Orthogonal equivariance allows
us to take $p=e_5$ and to diagonalize $A$ on $p^\perp$.  Write
\begin{equation}\label{cert:eigenvalues}
 A=\operatorname{diag}(x_1,x_2,x_3,x_4,\lambda),\qquad
 \lambda=-\sum_{i=1}^4x_i,
 \qquad R(x):=\sum_{i=1}^4x_i^2+\left(\sum_{i=1}^4x_i\right)^2=1.
\end{equation}
No sign or distinctness condition is imposed on the $x_i$; the normal
eigenvalue condition is exactly $\sum_i x_i>0$.
In these coordinates $(S_0)_{55}=0$ and
$(S_0)_{\alpha\alpha}=b_\alpha:=\lambda x_\alpha-x_\alpha^2$.
The only nonzero entries of $V$ are $V_{555}=1$ and the three permutations
of $V_{5\alpha\alpha}=b_\alpha$.  Since $C_{55jk}=0$, a direct component
formula for the boundary polynomial is therefore
\begin{equation}\label{cert:boundary-components}
 \mathcal S=-\lambda f+U_{55}+M_5+
       \sum_{\alpha=1}^4 b_\alpha
          (U_{\alpha\alpha}+C_{5\alpha5\alpha}
                                      +C_{5\alpha\alpha5}),
\end{equation}
where the fixed polynomials are evaluated at \eqref{cert:eigenvalues},
$p=e_5,g=0$.  This formula does not assume additional slot symmetries of
$U$ or $K$.

Equivalently, one may contract the boundary expression as a scalar graph
and then substitute
\begin{equation}\label{cert:boundary-substitution}
 p\cdot A^kp=\left(-\sum_i x_i\right)^k,\qquad
 \tr A^k=\sum_i x_i^k+\left(-\sum_i x_i\right)^k.
\end{equation}
For $k=0$, the latter is $5$ and the former is $1$.  The trace reductions
$\tr A=0$, $\tr A^2=1$ remain valid on $R=1$.  Thus either
\eqref{cert:boundary-components} or \eqref{cert:boundary-substitution}
evaluates the same scalar on all the allowed data.  Orthogonal invariance
under permutations of $p^\perp$ shows that $\mathcal S$ is symmetric in
$x_1,\ldots,x_4$.

Define the boundary remainder on this constraint set by
\[
 E_b=\Lambda+\frac14-\mathcal S.
\]
Expanding the two boundary Gram blocks and the Table~I expression gives a
polynomial representative with monomials of degree at most five.  A finite
way to check the degree assertion is to use
\eqref{cert:boundary-components}: discard positive $g$ powers, substitute
the diagonal matrix and normal vector in each of the $231$ rows, and add
the displayed terms.  All remaining monomials have degree at most five.
Each boundary $J_c$ has degrees one and two, so its squared terms have
degree at most four, or five after the $\texttt{H}$ weight.

The following convention gives a small exact list of its coefficients.
For an ordinary monomial $x^\alpha$ of degree $j\le5$, replace it by
\begin{equation}\label{cert:homogenization}
 x^\alpha R^{\lfloor(5-j)/2\rfloor}.
\end{equation}
This preserves its value on $R=1$ and gives degree four if $j$ is even and
degree five if $j$ is odd.  Perform the replacement before collecting
exponents.  Because $R$ is symmetric, it commutes with permutation averaging.
For a nondecreasing exponent vector $\beta$ put
\[
 m_\beta(x)=\operatorname{Av}(x^\beta).
\]
Expand \eqref{cert:homogenization} and sum the coefficients of all exponent
vectors whose nondecreasing rearrangement is $\beta$.  Their sum is the
coefficient of $m_\beta$ after averaging: indeed
$\operatorname{Av}(x^\alpha)=m_\beta$ for every member $\alpha$ of that
orbit, including when exponents repeat.  Thus sorting exponents and summing
coefficients introduces no extra factor of $24$.  Since $\mathcal S$ and
the constant are symmetric and $\Lambda$ was defined by averaging, the
resulting averaged polynomial still equals $E_b$ on $R=1$.

Precisely, using the row-pair expansion of Section~\ref{cert:collection}
for the two rank-$-1$ blocks, subtracting the boundary contributions from
Table~I with denominator $Q$, and adding the homogenized constant $1/4$,
one obtains
\begin{equation}\label{cert:boundary-expansion}
 E_b(x)=\sum_\beta e^{(b)}_\beta m_\beta(x)\quad\text{on }R=1,
\end{equation}
with exactly the following eleven nonzero coefficients:
\begin{equation}\label{cert:eleven-coefficients}
\begin{array}{c|r@{\qquad}c|r}
 \beta&Q^2 e^{(b)}_\beta&\beta&Q^2 e^{(b)}_\beta\\\hline
 (0,0,0,4)&-13479110 &(0,0,0,5)&17712712\\
 (0,0,1,3)&-93201124 &(0,0,1,4)&151038620\\
 (0,0,2,2)&-78543848 &(0,0,2,3)&244636976\\
 (0,1,1,2)&-139986162&(0,1,1,3)&310091962\\
 (0,1,2,2)&445177656 &(1,1,1,1)&-4861942\\
 (1,1,1,2)&157995170&&
\end{array}
\end{equation}
This coefficient computation consists solely of the explicitly stated
substitutions, multiplication of ordinary four-variable polynomials, and
coefficient addition.  It is also included in the accompanying exact
reproducer.

On $R=1$ every $|x_i|\le1$, hence $|m_\beta(x)|\le1$.  Adding the eleven
absolute integer coefficients gives
\[
 Q^2\sum_\beta|e^{(b)}_\beta|=1656725282.
\]
Since $\sum_i x_i=-\lambda>0$, \eqref{cert:boundary-gram} is nonnegative.
It follows from \eqref{cert:boundary-expansion} that
\[
 \mathcal S=\Lambda+\frac14-E_b
 \ge\frac14-\frac{1656725282}{62500000000}
 =\frac{6984137359}{31250000000}=\delta_*.
\]
This completes the proof of Proposition~\ref{cert:universal}.  The estimate
also persists at $\lambda=0$ by continuity; repeated tangent eigenvalues
cause no singularity in any of the formulas.

The pointwise comparisons and field bounds of Section~\ref{sec:flux}
are therefore established. Section~\ref{sec:cutoffs} now justifies their
integration across Hessian zeros and near the suspension tips.
